\documentclass[11pt]{article}
\usepackage[a4paper,margin=25mm]{geometry}
\usepackage{amsmath,amssymb,mathtools}
\usepackage[hidelinks]{hyperref}
\title{Complete massive triangle marginal at general Euclidean invariants}
\author{Phi3Lab isolated precision development}
\date{25 October 2026}
\begin{document}
\maketitle
The convention is mass squared one and one loop measure \(d^6k/\pi^3\).
Coupling and graph symmetry factors are not part of this scalar marginal.
Define
\[
 L(a,b,c)=\int_{x,y\geq0,\;x+y\leq1}\log F(x,y)\,dx\,dy,\qquad
 F=1+ax+by-ax^2-by^2+(c-a-b)xy .
\]
For \(a,b,c>0\), assume the strict Euclidean Gram condition
\(G=ab-(a+b-c)^2/4>0\).
The renormalized triangle is \(-L\). The simplex volume \(1/2\) is
already included in \(L\); it must not be multiplied a second time.

\section*{Boundary reduction}
Put \(d=(a+b-c)/2\), \(x_*=b(a-d)/(2G)\),
\(y_*=a(b-d)/(2G)\), and
\(\Delta=1+(ax_*+by_*)/2\).
The exact polynomial identity
\[
 (x-x_*)F_x+(y-y_*)F_y=2(F-\Delta)
\]
and the divergence theorem give
\[
 L=\frac{x_*E(b)+y_*E(a)+(1-x_*-y_*)E(c)}{2}
      -\frac12+\Delta J,\qquad
 J=\int_{\text{simplex}}\frac{dx\,dy}{F},
\]
where
\[
 E(t)=\int_0^1\log(1+t u(1-u))\,du
 =2\sqrt{\frac{t+4}{t}}\operatorname{arsinh}\frac{\sqrt t}{2}-2,
 \qquad E(0)=0.
\]
This is an identity for the full marginal, with no numerical quadrature
used in its closed evaluation.

\section*{Rationalized inverse triangle}
Let
\[
 \alpha=(a-b)^2-2c(a+b)+c^2=-4G,\quad
 \beta=2c(a+b-c),\quad r=\sqrt{c(c+4)}.
\]
For \(0\leq u\leq1\), define
\[
 R(u)^2=r^2+\beta u+\alpha u^2,\quad
 d(u)=c+u(a-b-c),\quad e(u)=c-u(a-b+c).
\]
Direct integration of the remaining simplex coordinate gives
\[
 J=\int_0^1\frac{du}{R(u)}
 \log\frac{(R+e)(R+d)}{(R-e)(R-d)} .
\]
Every logarithm argument is positive: exactly
\[
 R^2-d^2=4c[1+bu(1-u)],\qquad
 R^2-e^2=4c[1+au(1-u)] .
\]
With \(z=(R-r)/u\),
\[
 u=\frac{\beta-2rz}{z^2-\alpha},\qquad
 \frac{du}{R}=-\frac{2\,dz}{z^2-\alpha}.
\]
Write \(\kappa=\sqrt{-\alpha}\), \(z_0=\beta/(2r)\),
\(q_1=\sqrt{(a-b)^2+4c}\), and
\[
 \delta=z_1-z_0
 =\frac{4r^2\alpha-\beta^2}
        {2r(2r q_1+2r^2+\beta)}<0,\qquad z_1=q_1-r .
\]
This expression avoids subtracting nearly equal endpoint coordinates.

For \((\ell,s)=(a-b-c,b)\) or \((-a+b-c,a)\), and
\(\sigma\in\{1,-1\}\), factor
\[
 N_{\ell,\sigma}(z)
 =(r+\sigma c)(z^2-\alpha)+(\beta-2rz)(z+\sigma\ell).
\]
Its leading coefficient is \(\sigma c-r\), its middle coefficient
is \(\beta-2r\sigma\ell\), and its discriminant is
\(16c^2s(s+4)\).
The two roots \(\rho_{\ell,\sigma,\pm}\) are real and lie outside the
integration path. Set
\[
 B=\frac{\delta}{z_0-i\kappa},\qquad
 A_\rho=\frac{\delta}{z_0-\rho}.
\]
Then \(A_\rho\) is real and \(1+A_\rho>0\).

\section*{Dilogarithms and branches}
Define
\[
 {\cal P}(A,B)=\log(1-A/B)\log(1+B)
 -\operatorname{Li}_2\!\left(\frac{A(1+B)}{A-B}\right)
 +\operatorname{Li}_2\!\left(\frac{A}{A-B}\right).
\]
For real \(A\) with \(1+A>0\) and nonreal \(B\), this is
\[
 {\cal P}(A,B)=\int_0^1\frac{B\log(1+Av)}{1+Bv}\,dv .
\]
The branch justification is explicit. For
\(Z(v)=A(1+Bv)/(A-B)\),
\[
 \operatorname{Im}Z(v)
 =\frac{A\,\operatorname{Im}B\,(1+Av)}{|A-B|^2},\qquad
 (1-A/B)(1-Z(v))=1+Av>0 .
\]
The dilogarithm segment stays off its cut for nonzero \(A\), and the
principal logarithms combine without an untracked multiple of \(2\pi i\).
At \(A=0\), \({\cal P}=0\) by exact evaluation.
The result is
\[
 J=-\frac{2}{\kappa}\operatorname{Im}\left[
 2\log\frac{r+c}{r-c}\log(1+B)
 +\sum_{\ell,\sigma,\pm}
     \sigma\,{\cal P}(A_{\rho_{\ell,\sigma,\pm}},B)\right].
\]
The imaginary part implements the exact conjugate-pole identity; it is
not a rule for discarding a numerical error.

\section*{Degenerate limits and graph subtraction}
The exact case \(c=0,\ a=b\) is evaluated separately:
\(L(a,a,0)=E(a)/2\), including \(L(0,0,0)=0\).
The implementation rejects other degenerate or invalid inputs.
It does not clip a negative Gram determinant or replace a nearly
degenerate triple by the exact degenerate formula.

In a parent graph, a child's local \(T^0\) acts on all its boundary
momenta. After that child is renormalized, an overall \(T^0\) acts on
the physical external momenta and retains the host loop momenta.
These operations must not be interchanged or conflated.
For a unit external invariant and a single retained host momentum
of squared norm \(t\), the latter baseline is \(-E(t)/2\).

\section*{Reproducible validation and scope}
The proof job \texttt{triangle-complete-general-proof-01} completed
with exit zero under its independent supervisor.
Symbolic checks cover divergence, conic rationalization, all four
root discriminants, the endpoint identity, and the branch identity.
Nine scalar cases, including very small and very large invariants,
were compared with eighteen independent 180-digit integrals to
90 relative digits. Twenty-nine permutation checks, four unit-invariant
regressions, four exact degenerate cases and five invalid-domain
rejections passed.

The complete scalar identity is distinct from an acceptance of a
whole graph. Whole-graph results require their routed forest proof,
measure and integrability checks, interval implementation, independent
reference, fixed production policy and stored-request audit.
An empirical integration gate is not a rigorous quadrature-error
certificate.
\end{document}
